\clearpage
\chapter*{Fe de erratas del Exposé n\textsuperscript{o} 190}
\addcontentsline{toc}{chapter}{Fe de erratas del Exposé n\textsuperscript{o} 190}
\setcounter{footnote}{0}
\src{300}

\noindent\textbf{3. Técnica de descenso y teoremas de existencia en geometría
algebraica [TDTE I]. I: Generalidades. Descenso por morfismos fielmente planos.}

\begin{center}
[Séminaire Bourbaki, t. 12, 1959/60, n\textsuperscript{o} 190, 29 pp.]
\end{center}

\noindent\textbf{Página 1, línea 10.} --- Ahora parece excesivo decir que la
técnica de descenso está «en la base de la mayoría de los teoremas de
existencia en geometría algebraica». Esto es cierto en gran medida para las
técnicas no proyectivas que constituyen el objeto de los dos primeros exposés
de la serie TDTE I a VI, pero no para las técnicas proyectivas (TDTE IV, V,
VI).

\noindent\textbf{Página 5, línea 16.} --- No es necesario suponer que $\alpha$
sea un morfismo de $\mathcal F$-descenso.

\noindent\textbf{Página 20, observación.} --- Un morfismo
$S'\longrightarrow S$ cuasifinito, étale y sobreyectivo, o un morfismo
$\operatorname{Spec}(\widehat A)\longrightarrow\operatorname{Spec}(A)$, no es
siempre un morfismo de descenso estricto, aun si $A$ es un anillo local de una
curva algebraica sobre un cuerpo algebraicamente cerrado $k$ y
$S=\operatorname{Spec}(A)$. Así, se puede encontrar un morfismo propio y liso
$f:X\longrightarrow S$, que hace de $X$ un fibrado principal de base $S$ bajo
una curva elíptica $E$, tal que $f':X'\longrightarrow S'$ sea proyectivo, pero
$f$ no lo sea. Es, por tanto, al mismo tiempo un ejemplo de un fibrado
homogéneo principal no isotrivial bajo un esquema abeliano.

\noindent\textbf{Página 26, línea 9.} --- Léase «Chow--Lang» en lugar de
«Igusa--Lang».

Para diversos detalles acerca de la teoría del descenso, véanse SGA VI, VII y
VIII.
